% use input or include? % new page for each project or not?
\newif\ifdoinclude\doincludetrue

\ifdoinclude
  \newcommand{\meinclude}[1]{\include{#1}}
\else
  \newcommand{\meinclude}[1]{\input{#1}}
  \long\def\includeonly#1{\relax}
\fi

% If you want to see what package versions are used
% \listfiles % writes to logfile
  
% If only a part of the document is to be compiled:
% put \inlcudeonly{filelist} % in inlcudeonly.tex
\InputIfFileExists{includeonly.tex}{}{}
  
\areaset[10mm]{40pc}{675pt} % 

% Some simple packages
\usepackage{parskip}
\usepackage{paralist}
\usepackage{enumitem}
\usepackage{longtable}
\usepackage{blindtext}
\usepackage{phonenumbers}
\usepackage{amsmath,amssymb,amsthm}
\usepackage{eurosym}
\usepackage{csquotes}
\usepackage{pdfpages}
\usepackage[american]{isodate}
\usepackage{xspace}
\usepackage{changebar}
\usepackage{ulem}
\usepackage{booktabs}

% hyperref
\usepackage[colorlinks=false,
pdfstartview=FitV,
linkbordercolor=white,
% allcolors=black,
linkcolor=black,
citecolor=black,
citebordercolor=white,
pdfborderstyle={/S/U/W 1},
urlcolor=black,
urlbordercolor=white,
pagebackref=false]{hyperref}


% For showing defined labels:
% change final to draft in order to see the refs.
\usepackage[notref, notcite, final]{showkeys}
\renewcommand*\showkeyslabelformat[1]{%
                                      \rotatebox{45}{\normalfont
                                      \tiny
                                      \ttfamily#1}%
                                      }
\usepackage[disable, % comment to see todos.
obeyDraft,obeyFinal,textsize=small]{todonotes}

% xelatex specifics (better with ipa and stuff)
\usepackage{fontspec}
\usepackage{polyglossia}
\setdefaultlanguage[variant=american]{english}


% Bibliography and tweaks. Those tweaks aim at making the bibliography
% as short as possible which is usually what you desire.

% (it would be great to have different bibstlyes in the project
% descriptions, the cvs, and the main bibliographies. Unfortunately,
% "biblatex can't support multiple styles in the same document."
% (https://tex.stackexchange.com/questions/109234/change-biblatex-citation-style-after-preamble))
\usepackage[backend=biber,
	natbib=true,
	bibstyle=authoryear-icomp,
        firstinits=true,
        citestyle=authoryear-icomp, % supports ibid abbreviation.
	maxbibnames=5,
        maxcitenames=2,% as short as bearable?
        mincitenames=1,
	isbn=false,
	doi=false,
	eprint=true,
	url=false,
        autolang=hyphen,
        date=year% don't be impressed by the full date if given.
]{biblatex}
\IfFileExists{Concept.bib}{\addbibresource{Concept.bib}}{}

% Vor creating strings like V. Nachname and M. Mustermann":
% https://texwelt.de/wissen/fragen/5750/biblatex-funote-s-fehlt-wenn-f-oder-ff-nach-seitenzahl-folgt
\newrobustcmd*{\citefirstlastauthor}{\AtNextCite{\DeclareNameAlias{labelname}{first-last}}\citeauthor}

% to have titles like title: subtitle (the vikings: how they got where
% they went)
\renewcommand{\subtitlepunct}{: }

% removing some fields from the bibliography altogether and some
% conditional on the entry type:
\AtEveryBibitem{% Clean up the bibtex rather than editing it
 \clearlist{address}
 \clearfield{issn}
 \clearlist{location}
 \clearfield{month}
 \clearfield{series}
 \clearfield{pagetotal}
 
 \ifentrytype{article}{
  \clearlist{publisher}
  \clearname{editor}
 }{}
 \ifentrytype{book}{
   \clearfield{number}
 }{}
 \ifentrytype{inproceedings}{
   \clearfield{number}
 }{}
 \ifentrytype{incollection}{
   \clearfield{number}
 }{}
 \ifentrytype{collection}{
   \clearfield{number}
 }{}
}
% omit the pp. in bibliography.
\DeclareFieldFormat{pages}{#1}

% https://tex.stackexchange.com/questions/150461/replace-and-with-ampersand-in-bibliography-and-parenthetical-citations-using-b
\AtBeginBibliography{%
  \renewcommand*{\finalnamedelim}{%
    \ifnumgreater{\value{liststop}}{2}{%\finalandcomma
      % A, B, & C looks a bit stupid, so we change it to
      % A, B & C.
}{}%
    \addspace\&\space}%
}

%% https://tex.stackexchange.com/questions/22980/sentence-case-for-titles-in-biblatex
\DeclareFieldFormat{sentencecase}{\MakeSentenceCase{#1}}

% https://3d.bk.tudelft.nl/hledoux/blog/fiddling-biblatex/
% remove the quotes around the title of the paper
\DeclareFieldFormat[article, inbook, incollection, inproceedings, misc, thesis, unpublished]{title}{#1}

% https://tex.stackexchange.com/questions/10682/suppress-in-biblatex
\renewbibmacro{in:}{%
  \ifentrytype{article}{}{\printtext{\bibstring{in}\intitlepunct}}}
%\renewbibmacro{byeditor}{ed\addotspace} % wollte nicht.
%% Sentence case for title:subtitle in article and in* entries.
%% Titlecase else.
\renewbibmacro*{title}{%
  \ifthenelse{\iffieldundef{title}\AND\iffieldundef{subtitle}}
    {}
    {\ifthenelse{\ifentrytype{article}\OR\ifentrytype{inbook}%
      \OR\ifentrytype{incollection}\OR\ifentrytype{inproceedings}%
      \OR\ifentrytype{inreference}}
      {\printtext[title]{%
        \printfield[sentencecase]{title}%
        \setunit{\subtitlepunct}%
        \printfield[sentencecase]{subtitle}}}%
      {\printtext[title]{%
        \printfield[titlecase]{title}%
        \setunit{\subtitlepunct}%
        \printfield[titlecase]{subtitle}}}%
     \newunit}%
  \printfield{titleaddon}}
% end bibliography tweaks.

% for a consistent formatting of references between projects.
% https://tex.stackexchange.com/questions/64131/implementing-switch-cases
\usepackage{xifthen}
\newcommand{\project}[2]{\hypersetup{linkbordercolor=red}\hyperref[#1]{#1}}
\newcommand{\ifequals}[3]{\ifthenelse{\equal{#1}{#2}}{#3}{}}
\newcommand{\case}[2]{#1 #2} % Dummy, so \renewcommand has something to overwrite...
\newenvironment{switch}[1]{\renewcommand{\case}{\ifequals{#1}}}{}
\newcommand{\PP}[1]{%
    \begin{switch}{#1}%
        \case{A01}{\project{A01}{Stark}}%
        \case{B01}{\project{B01}{Baratheon \& Lannister}}%
        \case{MGK}{\project{MGK}{Stark \& Lannister}}%
        \case{INF}{\project{INF}{Stark, Baratheon \& Lannister}}%
        \case{Z01}{\project{Z01}{Stark \& Baratheon}}%
    \end{switch}%
  }
% For self references:  
\newcommand\PPs[1]{#1}
% Would be better to have only one command which detects if we have a
% self reference or not.

% Define your own commands here if you like:
\newcommand{\gethelp}[1]{\todo[inline, color=green]{#1}}
\normalem
\newcommand{\ul}[1]{\cbstart\uline{#1}\cbend\xspace}
\newcommand{\so}[1]{\cbstart\sout{#1}\cbend\xspace}

% for cost calculating tables
\newcommand{\ra}[1]{\renewcommand{\arraystretch}{#1}}

% language and font
\setmainfont{Latin Modern Sans}[
%                   SmallCapsFont={Latin Modern Roman Caps}
SmallCapsFont={[jkpssmsc8a.pfb]}
]
\setsansfont{Latin Modern Sans}
% following is so Junicode works on Mac; hopefully shouldn't
% break it for everyone else.
\newcommand{\Junicode}{\setmainfont
  [    Extension = .ttf,
  UprightFont = *,
  BoldFont = *-Bold,
  ItalicFont = *-Italic,
  BoldItalicFont = *-BoldItalic,
  ]{Junicode}}


%\overfullrule=5pt % to show overfull boxes

%%%% Titles and Other Strings %%%%
\newcommand{\cvpub}{Selected Publications}
\newcommand{\cvgeninf}{General Information}
\newcommand{\cveduc}{Academic Education with Qualifications}
\newcommand{\cvdeg}{Academic Degrees}
\newcommand{\cvdep}{Academic Career following End of Studies}
\newcommand{\cvmisc}{Miscellaneous}
\newcommand{\prbibhead}{Project-related publications by participating researchers}
\newcommand{\prbibheadsingle}{Project-related publications by participating researchers (all peer reviewed or monographs)}
\newcommand{\prsum}{Project description}
\newcommand{\resqu}{Summary and main research question}
\newcommand{\state}{Current state of research and preliminary work by principal investigator}
\newcommand{\statepl}{\state s}
\newcommand{\respr}{Research program and role within the proposed Collaborative Research Center}
\newcommand{\resprp}{Research program}
\newcommand{\role}{Role within the proposed CRC}
\newcommand{\pisg}{Principal investigator}
\newcommand{\pipl}{\pisg s}
\newcommand{\funds}{Project funding}
\newcommand{\peerlit}{Peer Reviewed Publications and Monographs}
\newcommand{\otherlit}{Other Publications}
\newcommand{\R}{\mathbb R}

% edit especially the following to your need
\newcommand{\crctitle}{CRC-Title}
\newcommand{\crcsubtitle}{CRC-Subtitle}
\newcommand{\thiscrc}{CRC \textit{Our CRC}}
\newcommand{\ABRA}{Uni~A}
\newcommand{\ABRB}{Uni~B}
\newcommand{\ABRC}{Uni~C}
\newcommand{\unia}{U Meinestadt\xspace}
\newcommand{\unialong}{Schön-Universität zu Meinestadt\xspace}
\newcommand{\uniaaddress}{Breite Straße~7, 22334 Meinestadt}
\newcommand{\uniclong}{Universität X}
\newcommand{\unicaddress}{Alte Straße 2, 98765 Bettenburg}
\newcommand{\uniblong}{Wissenschaftseinrichtung B\xspace}
\newcommand{\unibaddress}{Bäckerstr.~18, 33445 Deinestadt}
\newcommand{\instn}{Inst.\ f.\ Wissenschaft}

\newcommand{\wimi}{Research associate\xspace}
\newcommand{\sigsg}{Signature of the head of the applicant university}
\newcommand{\sigpl}{Signatures of the heads of the applicant universities}
% Those academic titles/positions work outside of names like in (\drphil)
\newcommand{\drphil}{Dr.~phil.\xspace}
\newcommand{\ma}{MA\xspace}
\newcommand{\dring}{Dr.-Ing.\xspace}
\newcommand{\dipling}{Dipl.-Ing.\xspace}
\newcommand{\phd}{PhD\xspace}
\newcommand{\pd}{PD\xspace}
\newcommand{\postdoc}{Post-doc\xspace}
% Those are always parts of academic titles as in \proflong Bertram.
% long version for headers, short for running text.
\newcommand{\proflong}{Prof.~} % Professor
\newcommand{\profshort}{Prof.~}
\newcommand{\uproflong}{Prof.~} % Universitätsprofessor
\newcommand{\uprofshort}{Prof.~}
\newcommand{\drlong}{Dr.~} % Doctor
\newcommand{\drshort}{Dr.~}
\newcommand{\drphillong}{Dr.~} % Doctor Philosophiae
\newcommand{\drphilshort}{Dr.~} % Dr. phil.
\newcommand{\drhclong}{Dr.~h.c.~} % Doctor h.c. or the like
\newcommand{\drhcshort}{Dr.~h.c.~}
\newcommand{\pdlong}{PD~} % Privatdozent
\newcommand{\pdshort}{PD~} % Priv.-Doz.
\newcommand{\phdlong}{Dr.~} % Ph.D.
\newcommand{\phdshort}{Dr.~} % Ph.D.
\newcommand{\diplmathlong}{Dipl.-Math.~} % Diplom-Mathematiker
\newcommand{\diplmathshort}{Dipl.-Math.~}
\newcommand{\diplinglong}{Dipl.-Ing.~} % Diplom-Mathematiker
\newcommand{\diplingshort}{Dipl.-Ing.~}
\newcommand{\diplinflong}{Dipl.-Inf.~} % Diplom-Informatiker
\newcommand{\diplinfshort}{Dipl.-Inf.~}


%%%% more specialized macros: %%%%

% subsections with numbers not occuring in toc:
\makeatletter
\newcommand{\notocsubsection}[1]{%
  \refstepcounter{subsection}%
  %\subsection*{\thesubsection \quad #1}
  \subsection*{\@arabic\c@subsection \quad #1}
}
\newcommand{\notocsubsubsection}[1]{%
  \refstepcounter{subsubsection}%
  %\subsection*{\thesubsection \quad #1}
  \subsection*{\@arabic\c@subsection.\@arabic\c@subsubsection \quad #1}
}
\makeatother

% customize bibs in cv & project descriptions:
\defbibheading{cvbib}[\bibname]{%
\subsubsection*{#1}}

\defbibheading{cvsubbib}[\bibname]{%
\paragraph*{#1}}

\defbibheading{prbib}[\bibname]{%
\notocsubsection{#1}}

\defbibheading{prsubbib}[\bibname]{%
% \paragraph*{#1}%
\textbf{#1}}

% explanations:
\newcommand{\erklaer}[1]{\begingroup\color{red}\begin{quote}#1\end{quote}\endgroup}
% switch off with uncommenting:
\renewcommand{\erklaer}[1]{}

% the same as environment, not as easily deactivated:
\newenvironment{erklaerung}{%
  \begingroup
  \color{red}
  \begin{quote}
  }
  {%
  \end{quote}
  \endgroup
}
%\renewenvironment{erklaerung}{}{}

% for time lines in the cv
\newlength\widest
\setlength\widest{5.5em}%
\newenvironment{cvtimeline}{%
  \begingroup
  % \begin{tabular}{p{.1\textwidth}p{.8\textwidth}}
  \begin{description}[leftmargin=\widest,
    labelwidth=\dimexpr\widest-0.5em\relax,
    font=\normalfont]
%  \begin[itemindent=1cm]{description}
  }
  {%
      %     \end{tabular}
  \end{description}
  \endgroup
}

% containing the pi-information in proj-desc
\newenvironment{giprenv}[1][.495]{%
  \begingroup
  \begin{minipage}[t]{#1\textwidth}
    \raggedright
  }
  {%
  \end{minipage}
  \endgroup
}
\newcommand{\gicolsep}{\hfill}
\newcommand{\girowsep}{
  
}

% formatting date ranges:
\newcommand{\fromto}[3][n]{%
  % 2016 2018: 2016 -- 2018
  % 2016 now: since 2016
  % 2016 2016: 2016
  % 2016 "": 2016
  % "" "" : ""
  % #1 j: The optional argument switches between brackets or no brackets.
\ifthenelse{\equal{#1}{j}\AND\NOT\(\isempty{#2}\AND\isempty{#3}\)}{(}{}%
\ifthenelse{\isempty{#3}\OR\equal{#2}{#3}}%
{#2}%
{\ifthenelse{\equal{#3}{now}}{since~#2}{#2--#3}}%
\ifthenelse{\equal{#1}{j}\AND\NOT\(\isempty{#2}\AND\isempty{#3}\)}{)}{}%
}



% for items within those timelines:
% format: \cvtlit{startyear}{endyear}{text}
\newcommand{\cvtlit}[3]{\item[\fromto{#1}{#2}]#3}

% for more general lists like the usual latex description environment:
\newenvironment{cvdesc}{%
  \begingroup
  \begin{description}[font=\normalfont]
  }
  {%
  \end{description}
  \endgroup
}
\newcommand{\cvitem}[2]{\item[#1] #2}

% bullet lists:
\newenvironment{cvitemize}{%
  \begingroup
  \begin{compactitem}
  }
  {%
  \end{compactitem}
  \endgroup
}
\newcommand{\cvit}{\item}
\newcommand{\cvitft}[3]{\item #1~\fromto[j]{#2}{#3}}

% If Argument contains
% ","
% " and "
% "\&"
% "/"
% \plurals return "s", else the empty string.
\newcommand{\plurals}[1]{\IfSubStr{#1}{,}{s}{\IfSubStr{#1}{ and }{s}{\IfSubStr{#1}{\&}{s}{\IfSubStr{#1}{/}{s}{}}}}}

% habilitation:
\newcommand{\cvhabil}[5][]{%[tit] field, univ, year, mentor
  \cvitem{Habilitation (#4)}{#2, #3,\ifthenelse{\NOT\isempty{#1}}{\textit{ #1\ifthenelse{\isempty{#5}}{.}{,}}}{} \ifthenelse{\isempty{#5}}{}{Mentor\plurals{#5}: #5}}
}

% field#1 degree#2
\newcommand{\fielddeg}[2]{
\ifthenelse{\isempty{#1}}{\ifthenelse{\isempty{#2}}{}{#2,}}{\ifthenelse{\isempty{#2}}{#1,}{#1 (#2)}}
}
% umhabilitation:
\newcommand{\cvumhabil}[3]{% field, univ, year
  \cvitem{Umhabilitation (#3)}{#1, #2}
}
% promotion.
%[tit] field, univ, year, supervisor, degree
\newcommand{\cvdoc}[6][]{
  \cvitem{Doctoral Degree (#4)}{\fielddeg{#2}{#6} #3,\ifthenelse{\NOT\isempty{#1}}{\textit{ #1\ifthenelse{\isempty{#5}}{.}{,}}}{} \ifthenelse{\isempty{#5}}{}{Supervisor\plurals{#5}: #5}}
}
% master
% [title]#1 field#2, univ#3, year#4, supervisor#5, degree#6
\newcommand{\cvmaster}[6][]{
  \cvitem{Master's Degree (#4)}% header + title
  {\fielddeg{#2}{#6} #3\ifthenelse{\NOT\isempty{#1}}{\textit{, #1}}{}\ifthenelse{\isempty{#5}}{}{, Supervisor\plurals{#5}: #5}
  }
} 

% [title] field, univ, year, supervisor, degree
\newcommand{\cvdipl}[6][]{
  \cvitem{Diplom (#4)}% header + title
  {\ifthenelse{\isempty{#2}}{\ifthenelse{\isempty{#6}}{}{#6,}}{\ifthenelse{\isempty{#6}}{#2,}{#2 (#6)}}\ifthenelse{\isempty{#3}}{}{, #3}\ifthenelse{\NOT\isempty{#1}}{\textit{, #1}}{}\ifthenelse{\isempty{#5}}{}{, Supervisor\plurals{#5}: #5}
  }
} 

% Staatsexamen
\newcommand{\cvseeins}[2][]{
  \cvitem{Erstes Staatsexamen (#2)}{#1\xspace(state examination)}
}
\newcommand{\cvsezwei}[2][]{
  \cvitem{Zweites Staatsexamen (#2)}{#1\xspace(state examination)}
}

% "project-related" in project-description
% not too unelegant, but it won't be possible (?) to refer by label to
% the generated section.
\newcommand{\prbib}[1]{
\begin{refsection}
\nocite{#1}
\printbibliography[heading=prbib,title={\prbibheadsingle}]
\end{refsection}
}


% peer reviewed bibliography (cv)
\newcommand{\peerbib}[1]{
\begin{refsection}
\nocite{#1}
\printbibliography[heading=cvbib,title={\peerlit}]
\end{refsection}
}

% peer reviewed bibliography (project description)
\newcommand{\prpeerbib}[1]{
  \begin{refsection}
    \nocite{#1}
    % \AtNextBibliography{\small} % might be forbidden
    \printbibliography[heading=prsubbib,title={\peerlit}]
  \end{refsection}
}

% other references (cv)
\newcommand{\otherbib}[1]{
  \begin{refsection}
    \nocite{#1}
    \printbibliography[heading=cvbib, title={\otherlit}]
  \end{refsection}
}

% other references (project descriptions)
\newcommand{\protherbib}[1]{
  \begin{refsection}
    \nocite{#1}
    % \AtNextBibliography{\small} % might be forbidden
    \printbibliography[heading=prsubbib, title={\otherlit}]
  \end{refsection}
}

% patents (cv):
\newcommand{\patentbib}[2]{%
  \ifthenelse{\isempty{#1}\AND\isempty{#2}}
  {}{
    \subsubsection*{c) Patents}}
\begin{refsection}
\nocite{#1}
\printbibliography[heading=cvsubbib, title={Filed}]
\end{refsection}
\begin{refsection}
\nocite{#2}
\printbibliography[heading=cvsubbib, title={Granted}]
\end{refsection}
}

\newcommand{\prbibs}[2]{
\notocsubsection{\prbibhead}
\ifthenelse{\isempty{#1}}{}{\prpeerbib{#1}}
\ifthenelse{\isempty{#2}}{}{\protherbib{#2}}}

% title of cv:
\newcommand{\cvhead}[1]{%
  \setcounter{secnumdepth}{0}%
  \setcounter{subsection}{0}% todo: das war auf einmal nötig?!
  % geht das eleganter?
  % Vielleicht einfach keine Nummern für die subsections?
  % subsubsection auch auf null?
  \markboth{#1}{#1}
  \thispagestyle{plain}
  \section{#1}}

\newcommand{\prhead}[3][]{%
  \setcounter{secnumdepth}{0}%
  \setcounter{subsection}{0}% todo: das war auf einmal nötig?!
  \setcounter{figure}{0}
  \section{#2: #3}\label{#2}
  \ifthenelse{\isempty{#1}}
  {\markboth{#2}{#2: #3}}
  {\markboth{#2}{#2: #1}} % shorttitle
  \thispagestyle{plain}
}


\makeatletter
% sections in cv:
\newcommand{\cvsection}[1]{%
  \refstepcounter{subsection}%
  \subsection*{\@arabic\c@subsection \quad #1}
}
% subsections in cv (not explicitely used):
\newcommand{\cvsubsection}[1]{%
  \refstepcounter{subsubsection}%
  \subsubsection*{\@arabic\c@subsection.\@arabic\c@subsubsection \quad #1}
}
\makeatother

% flx: todo: Was, wenn keine Homepage vorliegt (genauer)?
% flx: todo: Adresse wird schnell (viel) zu lang.
\makeatletter
% #1 [parental leave remark]
% #2 name
% #3 birth date
% #4 gender
% #5 institute+address
% #6 phone
% #7 mail
% #8 homepage
% #9 position
\newcommand{\ourphone}[1]{\phonenumber[foreign,foreign-area-code-sep=space]{#1}}
\let\textphone\ourphone
\newcommand*{\cvgi}{\begingroup\@makeother\_\@makeother~\@cvgi}
\newcommand*{\@cvgi}[9][]{#2\\
  \ifthenelse{\isempty{#3}}{}{*\printdate{#3},}%
  \ifthenelse{\equal{#4}{f}}{~female}{%
    \ifthenelse{\equal{#4}{m}}{~male}{~#4}
  }\\
  #5\\
  \ourphone{#6}\\
  \href{mailto:#7}{\texttt{#7}}\\
  \ifthenelse{\NOT\isempty{#8}}{
  \href{#8}{\texttt{#8}}\\}{}%
  #9\\
  #1
  \endgroup}
\makeatother

% typesetting the general information (project description)
\newcommand{\prgi}[7]{#1, *\printdate{#2}\\
  #3% inst 
%  #3\\ % inst 
\ifthenelse{\NOT\isempty{#4}}{, #4\\}{} % uni
#5\\ % address
% We don't want a newline here. If it's too wide, do something else about it.
\mbox{\phonenumber[foreign,foreign-area-code-sep=space]{#6} \textbullet{}
  \href{mailto:#7}{#7}}
}


\defbibheading{bibliographysection}{%
  \section*{#1}%
  \phantomsection
  \addcontentsline{toc}{chapter}{#1}
  \markboth{#1}{#1}%
  \thispagestyle{plain}
  \renewcommand*{\bibfont}{\small}%
}

% titelei
\subject{Proposed Collaborative Research Center}
\title{\crctitle}
\subtitle{\crcsubtitle}
\publishers{\unialong \\
  \uniblong}
\titlehead{\includegraphics[width=\textwidth]{misc/logos/logo1} \includegraphics{misc/logos/logo2}}

%  table of content
\setcounter{tocdepth}{3}

% numbering equations by section
\numberwithin{equation}{section}

\sloppy

